\documentclass[11pt]{article}
\usepackage[margin=1in]{geometry}
\usepackage{amsmath,amssymb,amsthm}
\usepackage{graphicx}
\usepackage{microtype}
\usepackage[hidelinks]{hyperref}

\newtheorem{theorem}{Theorem}
\newtheorem{lemma}{Lemma}
\newtheorem{corollary}{Corollary}
\newcommand{\Dir}{\operatorname{Dir}}
\newcommand{\Inv}{\operatorname{Inv}}

\title{Compact Operators Remain Compact\\
under the Bidual Direct-Limit Functor}
\author{Candidate solution packet for arXiv:2205.14385}
\date{21 July 2026}

\begin{document}
\maketitle

\noindent\textbf{Status.}
Candidate full solution, likely valid, pending human review.

\section{Source question}

Gwizdek constructs the tower
\[
 X_0=X,\qquad X_{n+1}=X_n^{**}
\]
with its canonical isometric embeddings, defines \(\Dir(X)\) as the
completion of the corresponding algebraic direct limit, and obtains
\(\Dir(T)\) by iterating biduals of a bounded operator \(T\)
\cite{source}. The compactness theorem in Section~2 assumes that \(X\) has
the approximation property. Immediately afterward, on PDF page~11, the
paper states:

\begin{quote}
``We conjecture that the approximation property assumption can be omitted.
\textbf{Hypothesis 2.14.} For any Banach space \(X\) if \(T\in B(X)\) is a
compact operator then \(\Dir(T)\) is a compact operator.''
\end{quote}

\begin{figure}[ht]
 \centering
 \includegraphics[width=\textwidth]{figures/open_problem_crop.png}
 \caption{Hypothesis 2.14 in the source paper, PDF page 11.}
\end{figure}

\section{Proof intuition}

The approximation property enters the source proof only to approximate a
compact operator by finite-rank operators. Instead, one can follow its compact
image through the bidual tower. If \(S:E\to F\) is compact and
\(K=\overline{S(B_E)}\), Goldstine's theorem and weak-star continuity give
\[
 S^{**}(B_{E^{**}})\subseteq \kappa_F(K).
\]
Thus bidualization does not create any new image points; it merely moves the
same norm-compact set into the canonical copy of the next space. At every
tower level, the image of the unit ball therefore represents a point of the
single compact set coming from level zero. Density then passes this containment
to the completed direct limit.

\section{The compact-image lemma}

For a Banach space \(E\), let
\(\kappa_E:E\to E^{**}\) be the canonical embedding and let \(B_E\) denote
the closed unit ball.

\begin{lemma}[compact images under bidualization]\label{lem:image}
Let \(E,F\) be Banach spaces and let \(S:E\to F\) be compact. If
\[
 K=\overline{S(B_E)},
\]
then
\[
 S^{**}(B_{E^{**}})\subseteq \kappa_F(K).
\]
In particular, \(S^{**}\) is compact.
\end{lemma}

\begin{proof}
Fix \(z\in B_{E^{**}}\). By Goldstine's theorem
\cite[Chapter~2]{megginson}, there is a net \((x_\alpha)\subseteq B_E\) such
that
\[
 \kappa_E x_\alpha\longrightarrow z
 \quad\text{in }\sigma(E^{**},E^*).
\]
The bidual map is weak-star--weak-star continuous, and naturality of the
canonical embeddings gives \(S^{**}\kappa_E=\kappa_F S\). Hence
\[
 \kappa_F(Sx_\alpha)=S^{**}(\kappa_E x_\alpha)
 \longrightarrow S^{**}z
 \quad\text{in }\sigma(F^{**},F^*).
\]
Since \(S\) is compact, \(K\) is norm compact. The set \(\kappa_F(K)\) is
therefore weak-star compact: the canonical embedding is norm-to-weak-star
continuous on \(K\). The weak-star topology is Hausdorff, so
\(\kappa_F(K)\) is weak-star closed. Every term of the displayed net belongs
to this set, and consequently its limit \(S^{**}z\) belongs to it as well.
The final assertion follows because \(S^{**}(B_{E^{**}})\) is contained in a
norm-compact set.
\end{proof}

\section{Main theorem}

For Banach spaces \(X,Y\), put
\[
 X_0=X,\quad X_{n+1}=X_n^{**},
 \qquad
 Y_0=Y,\quad Y_{n+1}=Y_n^{**}.
\]
Write \(\kappa^X_{n,n+1}\) and \(\kappa^Y_{n,n+1}\) for the canonical
embeddings, and write
\[
 \iota_n^X:X_n\to\Dir(X),
 \qquad
 \iota_n^Y:Y_n\to\Dir(Y)
\]
for the canonical isometries. Given \(T:X\to Y\), set
\(T_0=T\) and \(T_{n+1}=T_n^{**}\). The definitions give
\begin{equation}\label{eq:naturality}
 \Dir(T)\iota_n^X=\iota_n^Y T_n,
 \qquad
 \iota_{n+1}^Y\kappa^Y_{n,n+1}=\iota_n^Y.
\end{equation}

\begin{theorem}\label{thm:main}
Let \(X,Y\) be arbitrary Banach spaces. If \(T:X\to Y\) is compact, then
\[
 \Dir(T):\Dir(X)\longrightarrow\Dir(Y)
\]
is compact.
\end{theorem}

\begin{proof}
For every \(n\geq0\), define
\[
 K_n=\overline{T_n(B_{X_n})}\subseteq Y_n.
\]
The set \(K_0\) is compact. Lemma~\ref{lem:image}, applied first to \(T_0\)
and then inductively to \(T_n\), shows that every \(T_n\) is compact and
\begin{equation}\label{eq:nesting}
 K_{n+1}\subseteq \kappa^Y_{n,n+1}(K_n)
 \qquad(n\geq0).
\end{equation}
Indeed, the lemma gives the inclusion before closure, and its right-hand side
is compact and hence norm closed.

Combining \eqref{eq:nesting} with \eqref{eq:naturality} and iterating gives
\begin{equation}\label{eq:fixedK}
 \iota_n^Y(K_n)\subseteq \iota_0^Y(K_0)
 \qquad(n\geq0).
\end{equation}
Let
\[
 A_X=\bigcup_{n\geq0}\iota_n^X(X_n).
\]
This algebraic direct limit is dense in \(\Dir(X)\). If
\(a\in A_X\cap B_{\Dir(X)}\), then \(a=\iota_n^X x\) for some
\(x\in B_{X_n}\), since \(\iota_n^X\) is an isometry. Equations
\eqref{eq:naturality} and \eqref{eq:fixedK} therefore imply
\begin{equation}\label{eq:denseball}
 \Dir(T)(A_X\cap B_{\Dir(X)})\subseteq \iota_0^Y(K_0).
\end{equation}

The set \(A_X\cap B_{\Dir(X)}\) is dense in \(B_{\Dir(X)}\). To see this,
approximate any \(a\in B_{\Dir(X)}\) by \(a_j\in A_X\), and replace
\(a_j\) with
\[
 \frac{a_j}{\max\{1,\lVert a_j\rVert\}}.
\]
These normalized approximants still converge to \(a\). Finally,
\(\iota_0^Y(K_0)\) is norm compact and hence closed. By continuity of
\(\Dir(T)\), the containment \eqref{eq:denseball} passes to the completed
unit ball:
\[
 \Dir(T)(B_{\Dir(X)})\subseteq \iota_0^Y(K_0).
\]
Thus the image of the unit ball is relatively compact, proving that
\(\Dir(T)\) is compact.
\end{proof}

\begin{corollary}
Hypothesis 2.14 of \cite{source} holds without the approximation property. If
\(T\in B(X)\) is compact, then both \(\Dir(T)\) and \(\Inv(T)\) are compact.
\end{corollary}

\begin{proof}
Theorem~\ref{thm:main} gives the assertion for \(\Dir(T)\). Theorem~2.9 of
the source identifies \(\Inv(T)\) with the adjoint of \(\Dir(T)\). Schauder's
theorem \cite[Chapter~3]{megginson} says that the adjoint of a compact
operator is compact.
\end{proof}

\section{Verification and limitations}

The adversarial audit checks the unit-ball form of Goldstine's theorem, the
weak-star closedness of \(\kappa_F(K)\), the direction of the inclusions in
\eqref{eq:nesting}, and density of the algebraic direct-limit unit ball. No
gap was found. No computation is relevant to this qualitative argument.

A bounded novelty search on 21 July 2026 covered the run's lightweight
indexes, arXiv/web queries for the exact conjecture sentence, paper title,
arXiv id, author, Hypothesis~2.14, and close keyword variants, as well as
visible later-title/citation results. No explicit later resolution was found.
A 2025 Polish dissertation by the same author visible in the search results
still describes the compactness theorem with an extra assumption and traces it
to the source paper. This is negative search evidence, not an exhaustive
review or a priority claim.

The result is stronger than the stated hypothesis because it allows compact
maps \(X\to Y\). Status remains \emph{candidate full solution, likely valid}
until human review. The recommended review is to recheck
Lemma~\ref{lem:image} and the completion step in
\eqref{eq:denseball}.

\begin{thebibliography}{9}

\bibitem{source}
S.~Gwizdek,
\emph{Isometric direct limits of bidual Banach spaces},
arXiv:2205.14385 (2022; source PDF version dated 2023).

\bibitem{megginson}
R.~E.~Megginson,
\emph{An Introduction to Banach Space Theory},
Graduate Texts in Mathematics 183, Springer, 1998.

\end{thebibliography}

\end{document}
